% !TeX program = xelatex
\documentclass[12pt,a4paper]{article}
\usepackage[margin=2cm]{geometry}
\usepackage{fontspec}
\usepackage{amsmath,amssymb}
\usepackage{enumitem}
\setmainfont{Times New Roman}
\setlength{\parindent}{0pt}
\setlength{\parskip}{6pt}

\begin{document}
\begin{center}
  {\Large\bfseries ĐỀ THI TUYỂN SINH VÀO LỚP 10 THPT}\\
  {\bfseries Thành phố Hà Nội, năm học 2008--2009}\\
  {\itshape Phần Đại số}
\end{center}

\textbf{Câu 1.} Cho biểu thức
\[
P=\left(\frac{1}{\sqrt{x}}+\frac{\sqrt{x}}{\sqrt{x}+1}\right)
:\frac{\sqrt{x}}{x+\sqrt{x}}.
\]
\begin{enumerate}[label=\alph*)]
  \item Rút gọn $P$.
  \item Tính giá trị của $P$ khi $x=4$.
  \item Tìm giá trị của $x$ để $P=\dfrac{13}{3}$.
\end{enumerate}

\textbf{Câu 2.} Giải bài toán sau bằng cách lập phương trình hoặc hệ phương trình:

Tháng thứ nhất hai tổ sản xuất được 900 chi tiết máy. Tháng thứ hai, tổ I vượt mức 15\% và tổ II vượt mức 10\% so với tháng thứ nhất, vì vậy hai tổ sản xuất được 1010 chi tiết máy. Hỏi tháng thứ nhất mỗi tổ sản xuất được bao nhiêu chi tiết máy?

\textbf{Câu 3.} Cho parabol $(P):y=\dfrac{x^2}{4}$ và đường thẳng $(d):y=mx+1$.
\begin{enumerate}[label=\alph*)]
  \item Chứng minh rằng với mọi giá trị của $m$, đường thẳng $(d)$ luôn cắt parabol $(P)$ tại hai điểm phân biệt $A$ và $B$.
  \item Tính diện tích tam giác $AOB$ theo $m$, với $O$ là gốc tọa độ.
\end{enumerate}

\textbf{Câu 5.} Tìm giá trị nhỏ nhất của biểu thức
\[
Q=(x-1)^4+(x-3)^4+6(x-1)^2(x-3)^2.
\]

\end{document}
